% =====================================================================
\section{Windowed register file}
\label{sec:regfile}
% =====================================================================

The 64 physical registers P0--P63 form 16 blocks of 4 registers (B0--B15).
The program sees 8 of them as R0--R7. Because the window moves by 4
registers, adjacent windows share 4 registers, which are used to pass
arguments and results between functions.

\subsection{Mapping}

Logical register $R_n$ ($n = 0 \ldots 7$) in window CWP maps to the physical
register
\begin{equation}
  P_{\mathit{idx}} = (4 \cdot \mathrm{CWP} + n) \bmod 64
  \label{eq:window}
\end{equation}

Example: with CWP\,=\,3, registers R0--R3 are P12--P15 (block B3) and R4--R7
are P16--P19 (block B4). After \texttt{WINC} (CWP\,=\,4), the former R4--R7
become R0--R3 (Figure~\ref{fig:windows}).

\begin{figure}[htbp]
  \centering
  % ---------------------------------------------------------------------
%  Register windows: 16 blocks of 4 registers, window = 2 blocks
%  1 unit = 0.2 mm, y grows downwards. Block i starts at x = 24 + 44*i.
% ---------------------------------------------------------------------
\begin{tikzpicture}[
    x=0.2mm, y=-0.2mm,
    >={Stealth[length=1.8mm,width=1.4mm]},
    conn/.style={draw=lineclr, line width=0.6pt},
    lbl/.style={font=\scriptsize},
  ]
  % --- 16 blocks ------------------------------------------------------
  \foreach \i in {0,...,15} {
    \pgfmathsetmacro{\bx}{24+44*\i}
    \ifnum\i=4
      \draw[draw=accent, fill=accent!12, line width=1.1pt, rounded corners=1mm]
        (\bx,96) rectangle ++(40,40);
      \node[font=\scriptsize\bfseries] at (\bx+20,116) {B\i};
    \else
      \draw[draw=lineclr, line width=0.6pt, rounded corners=1mm]
        (\bx,96) rectangle ++(40,40);
      \node[font=\scriptsize] at (\bx+20,116) {B\i};
    \fi
  }

  % --- window CWP = 3 (blocks 3, 4) -----------------------------------
  \draw[conn] (156,90) -- (156,80) -- (240,80) -- (240,90);
  \node[lbl, anchor=south west] at (156,76) {CWP = 3: R0--R3 = B3, R4--R7 = B4};

  % --- window CWP = 4 (blocks 4, 5) -----------------------------------
  \draw[conn] (200,142) -- (200,152) -- (284,152) -- (284,142);
  \node[lbl, anchor=north west] at (200,156) {CWP = 4 (after \texttt{WINC}): R0--R3 = B4, R4--R7 = B5};

  % --- WSTART marker --------------------------------------------------
  \node[lbl, text=black!60, anchor=south] at (44,92) {WSTART};

  % --- circular wrap-around -------------------------------------------
  \draw[conn,->] (704,136) -- (704,200) -- (44,200) -- (44,136);
  \node[lbl, text=black!60, anchor=north] at (374,204)
    {B15 $\rightarrow$ B0: circular buffer, index mod 16};

  % --- notes ----------------------------------------------------------
  \node[font=\footnotesize, anchor=west] at (24,248)
    {B4 is shared: caller's R4--R7 = callee's R0--R3 (arguments and results)};
  \node[font=\footnotesize, text=black!60, anchor=west] at (24,272)
    {\texttt{WINC} with (CWP + 2) mod 16 = WSTART $\rightarrow$ WOVF exception, oldest block is spilled};
\end{tikzpicture}

  \caption{Register windows: 16 blocks of 4 registers, adjacent windows share one block}
  \label{fig:windows}
\end{figure}

Each window occupies two consecutive blocks, so \texttt{WINC} allocates one new
block while the other stays shared with the caller's window.

\subsection{Calling convention}

\begin{enumerate}
  \item The caller places up to 4 arguments in R4--R7 and executes
        \texttt{CALLW} (= \texttt{WINC} + \texttt{CALL}).
  \item The callee sees the arguments in R0--R3; R4--R7 are its working
        registers and, at the same time, the argument area for further calls.
  \item The callee returns results in R0--R3 and executes \texttt{RETW}
        (= \texttt{WDEC} + \texttt{RET}); the caller reads them in R4--R7.
  \item More than 4 arguments are passed on the stack.
\end{enumerate}

\subsection{Requirements}

\begin{reqtable}
\req{REG-001}{The processor \MUST{} have 64 general-purpose 8-bit registers
  (P0--P63).}
\req{REG-002}{At any time, 8 registers R0--R7 \MUST{} be visible, mapped
  according to Equation~\eqref{eq:window}.}
\req{REG-003}{The window pointer CWP \MUST{} be 4 bits wide (16 positions) and
  \MUST{} be updated modulo 16.}
\req{REG-004}{The \texttt{WINC} instruction \MUST{} increment CWP by 1 (the
  window moves up by 4 registers); \texttt{WDEC} \MUST{} decrement it by 1.}
\req{REG-005}{The register file \MUST{} support two reads and one write in a
  single cycle.}
\req{REG-006}{Register pairs (R0:R1, R2:R3, R4:R5, R6:R7) \MUST{} be usable as
  16-bit address pointers, with the low byte in the even register.}
\req{REG-007}{The register WSTART (4 bits) \MUST{} point to the oldest window
  whose contents are still held in the register file.}
\req{REG-008}{Executing \texttt{WINC} when $(\mathrm{CWP} + 2) \bmod 16 =
  \mathrm{WSTART}$ \MUST{} raise the window overflow exception (WOVF) before
  CWP is changed. This allows up to 15 simultaneously live windows (14 nested
  calls) without spilling to memory.}
\req{REG-009}{Executing \texttt{WDEC} when $\mathrm{CWP} = \mathrm{WSTART}$
  \MUST{} raise the window underflow exception (WUNF) before CWP is changed.}
\req{REG-010}{After a WOVF/WUNF handler returns (\texttt{RETI}), the
  instruction that raised the exception \MUST{} be executed again.}
\req{REG-011}{The exception handler \MUST{} be able to access any physical
  register through the SFRs PRIDX (index 0--63) and PRDAT (data), in order to
  save the WSTART block to the stack (spill) or restore it (fill).}
\req{REG-012}{Register contents after reset \MAY{} be undefined; the
  implementation \SHOULD{} clear them in simulation to simplify debugging.}
\end{reqtable}
